\paragraph*{\texttt{valueShift} 的适用范围}

设 \(P\) 为素数，\(v_i=f(i)\ (0\le i\le n)\)、\(\deg f\le n\)，且输入
\(v_i\) 已规范到 \([0,P)\)。函数
\[
\texttt{valueShift(v,c)}
\]
返回
\[
f(c),f(c+1),\ldots,f(c+n).
\]
要求 \(v\) 非空，调用前已将 \(c\) 规范到 \(0\le c<P\)，并满足
\[
2n+1<P,\qquad n<c<P-n.
\]
等价地，已知点集 \(\{0,\ldots,n\}\) 与目标点集
\(\{c,\ldots,c+n\}\) 在 \(\mathbb F_P\) 中不相交。当前实现不支持两个
点集重叠的平移；此时 Lagrange 公式会出现需要先消去的零因子，不能直接
计算模逆。时间复杂度为 \(O(M(n))\)，空间复杂度为 \(O(n)\)。
